% ---------------------------------------------------------------------------------------------------------
% 考察・含意・限界・まとめ
% ---------------------------------------------------------------------------------------------------------
\section{考察}
\label{sec:discussion}
\textbf{双対ラベルは何を改善し，何を改善しなかったか．}
双対ラベルは，価値の勾配の精度（未知の体格での順位相関 \fitSobHoldRho 対 \fitValHoldRho），
0.8\,s 先の候補の順位付け（最良を選ぶ割合 \rankSobLong\,\%対 \rankValLong\,\%），
閉ループで離手に至る割合（\clSmpcRel\,\%対 \clVmpcRel\,\%）を改善した．
しかし成功率は改善しなかった．改善は，成功に必要な段階のうち最初の一つにとどまっている．
成功には，離手の瞬間の状態が，共通の着地反射で捕捉できる範囲に入っている必要がある．
参照解を 2\,ms の制御周期で再生しても成功率が \clReplaySuccInt\,\%であること，
参照解の状態を追従する MPC が 1 本も捕捉できないことは，この範囲が非常に狭いことを示す．
学習した場による MPC は，その範囲に入る前に，振りの途中で錐，壁，期限のいずれかで失敗している．

\textbf{なぜ場の精度が足りないか．}
MPC の地平 0.3\,s では，候補の間の評価値の差（中央値 \spreadShortMed）は場の誤差よりはるかに小さく，
オラクルの解き直しのばらつきと同程度であった（第~\ref{sec:ranking}節）．
地平を 0.8\,s にすれば候補は区別できるが，その地平の MPC でも成功率は変わらなかった．
候補を正しく順位付けできることは，振りの全体を通して制約の内側にとどまることを保証しない．
さらに，価値は最適解の近くで平坦であり（第~\ref{sec:audit}節），costate のラベルは方向は信頼できても大きさに 3 割の誤差がある．
教師の側の精度にも限界がある．

\textbf{計画の側の問題．}
最適解は可動域の端，壁，錐の境界の上を通る（図~\ref{fig:timeseries}）．
計画のメッシュを細かくすると再生は改善し，離手時刻を 1\,ms ずらすと悪化する（第~\ref{sec:replay}節）．
一方，制約を一律に内側へ締める余裕は成功率を上げず，離手時刻の誤差を織り込んだ計画は解けなかった（第~\ref{sec:margins}節）．
実行に耐える計画を作るには，一律の余裕ではなく，実行の誤差がどの制約をどれだけ破るかを織り込んだ定式化が要る．

\textbf{データの集め方．}
制御器が訪れた状態を 100 回ほど解き直して加えるだけで，離手率は同じデータ量の参照近傍だけの場合の 5 倍近くになった（第~\ref{sec:data}節）．
参照解の近くのデータをいくら増やしても，制御器が実際に行き着く状態は覆えない．
この方向（訪問状態の解き直しの反復）は，本稿では 1 回しか行っていない．

\section{ロボティクスへの含意}
\label{sec:robotics}
本研究の対象は人体だが，方法はロボットにそのまま移る．
第一に，必要容量の地図は設計仕様である．3\,cm の縁を両手で捕捉する動作に要る把持力（体重の \capBWmin〜\capBWmax 倍）と，
関節トルクを増やしても把持力の要求がほとんど下がらないという結果は，懸垂移動ロボットではアクチュエータよりもグリッパが律速になることを意味する．
第二に，最適化器と環境が同じ記号力学を共有する設計は，学習した制御器を厳密な最適解と同じ土俵で測るための実験基盤になる．
計画の再生が失敗する原因を，制御周期，離手時刻，メッシュ，開始相に分解できたのはこの設計による．
第三に，包絡線定理による転移は，アクチュエータ容量の $\pm5$\,\%程度の違いなら解き直さずに最適値を予測できることを示した．
個体差のあるロボットへの一次の補正として使える．

\section{限界と今後}
\label{sec:limits}
平面モデルは矢状面だけを扱い，体の向きの反転を自由としている．ひねりに要る角運動量と，片手ずつの捕捉の負荷は 3 次元モデルの課題である．
把持容量は両手の合計の力で表し，指の疲労，皮膚の摩擦，個人差を含まない．
得られた必要容量は局所解である可能性があり，上界として読むべきである．
捕捉直後の把持利用率は環境の積分刻みについて収束しておらず，閉ループの $\Upk$ の絶対値には接触モデルの不確かさが残る．
閉ループの比較は 1 つの装置周期と 4 体格に限られ，学習の反復は 3 回である．
成功率が 0 の制御器どうしの比較では，離手率のような途中の指標の差しか検出できない．
今後は，(i)~実行の誤差を織り込んだ計画（離手時刻と状態の誤差に対する頑健化），
(ii)~訪問状態の解き直しの反復，(iii)~着地反射を学習の対象に含めること，(iv)~3 次元モデルへの拡張を進める．

\section{まとめ}
クリフディメンションの背面飛び移りを，装置運動，3\,cm の突起，有限の把持容量，前腕の干渉を含む力学課題として定式化し，
\gridN 条件の軌道最適化から必要把持容量を求めた．
必要容量は体重の \capBWmin〜\capBWmax 倍で，装置周期にほとんど依らず，捕捉後の保持が決めている．
最適化の双対変数を学習する双対場学習は，価値の勾配，候補の順位付け，離手率を改善したが，
成功率は双対ラベルの有無にかかわらず \clSmpcSuccInt\,\%であった．
参照解の再生でも成功率は \clReplaySuccInt\,\%で，この課題の難しさは，制約の境界の上を通る最適解を実行することにある．
コードと結果は \url{https://github.com/omotes24/CliffDimension} にある．
